\documentclass[11pt]{article}
\usepackage[a4paper,margin=26mm]{geometry}
\usepackage[T1]{fontenc}
\usepackage{lmodern,amsmath,amssymb,amsthm,mathtools,booktabs}
\usepackage[colorlinks=true,linkcolor=blue,urlcolor=blue,citecolor=blue]{hyperref}
\newtheorem{theorem}{Theorem}
\newtheorem{lemma}[theorem]{Lemma}
\newtheorem{corollary}[theorem]{Corollary}
\newcommand{\tr}{\operatorname{tr}}
\newcommand{\Cov}{\operatorname{Cov}}
\newcommand{\Del}{\mathcal D}
\title{Weighted Forest Rayleigh Inequalities on Every Wheel\\\large A two-by-two matrix certificate and its consequences}
\author{Research note}
\date{14 November 2026}
\begin{document}
\maketitle
\begin{abstract}
We give a self-contained proof that the forest polynomial of every wheel is Rayleigh at every positive assignment of edge weights. The argument uses an exact transfer formula, inversion of rim weights, and elementary invariants of a semigroup of two-by-two matrices. Three explicit certificates treat the three types of edge pairs. A sharp relaxation parameter also gives the independent-set Rayleigh property for all whirls. Consequences include minor and two-sum closure, conditional pairwise negative correlation, negative cylinder inequalities, concentration bounds, and the corresponding weighted connected-subgraph inequality on wheels. This is a proof draft for infinite families, not a resolution of the forest conjecture for arbitrary graphs. External review and a comprehensive priority check remain pending.
\end{abstract}

\section{Statement and relation to existing work}
Let $W_n=K_1+C_n$, $n\ge3$, have hub $h$ and cyclic rim vertices $v_1,\ldots,v_n$. Write $s_i>0$ for the weight of $hv_i$ and $r_i>0$ for the weight of $v_iv_{i+1}$. For any graph $G$, its unrooted spanning-forest polynomial is
\[
 Z_G(\boldsymbol w)=\sum_{A\subseteq E(G):\,A\text{ acyclic}}\prod_{e\in A}w_e.
\]
Isolated vertices are allowed. Define $\Del_{ef} Z=Z_eZ_f-ZZ_{ef}$ for distinct edges, where subscripts denote partial derivatives.

\begin{theorem}\label{thm:main}
For every $n\ge3$, every positive assignment of the $2n$ edge weights, and all distinct edges $e,f$ of $W_n$,
\[
\Del_{ef}Z_{W_n}>0.
\]
Consequently, under the probability law $\Pr(A)=\prod_{e\in A}w_e/Z_{W_n}$ on forests, every two distinct edge indicators are strictly negatively correlated.
\end{theorem}

Indeed, $\Cov(X_e,X_f)=-w_ew_f\Del_{ef}Z/Z^2$. Strictness is asserted for the wheel with all weights positive; its minors satisfy the non-strict inequality.

Bado's 2026 paper proves the unit-weight case for all wheels. Its Remarks 7.1--7.2 distinguish this from, and propose studying, the full weighted case \cite{Bado}. The transfer formula below is equivalent to that paper's multivariate formula. Our contribution in this note is the positivity argument with arbitrary weights, rather than the transfer formula itself. The larger arbitrary-graph problem remains open in the recent literature \cite{CR}.

This investigation was motivated by VMRC [016] (rooted forest determinants) and [160] (an obstruction to a prescribed forest sum-of-squares ansatz) \cite{VMRC}. Neither manuscript is assumed as a theorem in the proof below. In particular, rooted and unrooted forest measures are not identified.

\section{Exact transfer formula}
Put
\[
 P(s)=\begin{pmatrix}1&0\\s&1\end{pmatrix},\qquad
 R(t)=\begin{pmatrix}1+t&t\\0&1\end{pmatrix},\qquad t_i=r_i^{-1},
\]
and
\[
 H(\boldsymbol s,\boldsymbol t)=\tr\!\left(P(s_1)R(t_1)\cdots P(s_n)R(t_n)\right)-2.
\]
\begin{lemma}\label{lem:transfer}
The forest polynomial is
\begin{equation}\label{eq:transfer}
 Z_{W_n}(\boldsymbol s,\boldsymbol r)=\left(\prod_{i=1}^n r_i\right)
 H(\boldsymbol s,\boldsymbol r^{-1}).
\end{equation}
\end{lemma}
\begin{proof}
Write $R(t)=I+tU$, where $U=uv^{\mathsf T}$, $u=(1,0)^{\mathsf T}$ and $v=(1,1)^{\mathsf T}$. Expand by the subset of selected $t_iU$ factors. Selecting none contributes $\tr P(\sum_i s_i)=2$, which is cancelled.

A nonempty selected subset specifies the missing rim edges. Its complementary rim subgraph consists of paths, including singleton paths. The rank-one factors split the cyclic trace into one factor per path component $C$. Since $P(a)P(b)=P(a+b)$, this factor is
\[
 v^{\mathsf T}P\!\left(\sum_{i\in C}s_i\right)u=1+\sum_{i\in C}s_i.
\]
A rim path component can contain at most one selected spoke: two would create a cycle through the hub. Thus this factor counts precisely the valid spoke choices in that component. Multiplication by $\prod r_i$ changes the missing-edge weight $\prod_{i\in A}t_i$ into $\prod_{i\notin A}r_i$. This counts every forest once.
\end{proof}

For comparison, setting $Q(r)=rR(1/r)$ gives $Z=\tr\prod P(s_i)Q(r_i)-2\prod r_i$. Conjugation by $C=\left(\begin{smallmatrix}1&1\\1&0\end{smallmatrix}\right)$ yields
\[
 CP(s)Q(r)C^{-1}=\begin{pmatrix}1+s+r&sr\\1&r\end{pmatrix},
\]
the transfer matrix in \cite{Bado}.

\section{Two semigroup invariants}
Let $\mathcal S$ be the semigroup of all finite products of $P(s)$ and $R(t)$ with $s,t\ge0$, including the empty product $I$. Every member has nonnegative entries and determinant at least one.

\begin{lemma}\label{lem:invariants}
For $X=\left(\begin{smallmatrix}a&b\\c&d\end{smallmatrix}\right)\in\mathcal S$, let $D_X=\det X$. Then
\begin{equation}\label{eq:uv}
 D_X(a-b)\ge a,\qquad D_Xd\ge b+d.
\end{equation}
With $S=\left(\begin{smallmatrix}1&1\\0&-1\end{smallmatrix}\right)=S^{-1}$ and
\[
 X'=SXS=\begin{pmatrix}\alpha&\beta\\\gamma&\delta\end{pmatrix},
\]
one has
\begin{equation}\label{eq:conj}
 \alpha\ge D_X,\qquad \beta\ge0,\qquad\gamma\le0,
 \qquad \alpha\delta\le D_X.
\end{equation}
\end{lemma}
\begin{proof}
For the first inequality define $I(X)=D_X(a-b)-a$. Right multiplication by $P(s)$ or $R(t)$ gives, respectively,
\[
 I(XP(s))=I(X)+s b(D_X-1),\qquad I(XR(t))=(1+t)I(X).
\]
Both preserve nonnegativity, and $I(I)=0$.

For the second define $J(X)=D_Xd-b-d$. Left multiplication gives
\[
 J(P(s)X)=J(X)+s b(D_X-1),\qquad J(R(t)X)=(1+t)J(X).
\]
Building the word from the opposite end proves $J(X)\ge0$.

For the conjugated assertions, write $A=a+c$ and $B=b+d$, the column sums of the original matrix. Initially $A=B=D_X=1$. On right multiplication by $P(s)$,
\[
 (A,B,D_X)\longmapsto(A+sB,B,D_X);
\]
on right multiplication by $R(t)$,
\[
 (A,B,D_X)\longmapsto((1+t)A,tA+B,(1+t)D_X).
\]
These transformations preserve $A\ge B$ and $A\ge D_X$. Direct conjugation gives $\alpha=A$, $\beta=A-B$, $\gamma=-c$. Finally, $\alpha\delta=D_X+\beta\gamma\le D_X$.
\end{proof}

\section{Three certificates}
Each variable occurs in exactly one affine matrix factor. Hence $H$ is multiaffine. For any multiaffine expression, its Rayleigh difference in two variables is independent of those variables. We can therefore set the marked variables to zero and replace their two factors by $I$. The remaining factors form two cyclic segments $X,Y\in\mathcal S$.

\subsection{Two spokes}
Let $E=\left(\begin{smallmatrix}0&0\\1&0\end{smallmatrix}\right)$. The two marked factors are $I+sE$ and $I+zE$. At $s=z=0$, cyclic differentiation gives
\[
 H_s=(XY)_{12},\qquad H_z=(YX)_{12},\qquad H_{sz}=X_{12}Y_{12}.
\]
For arbitrary two-by-two matrices,
\[
 (XY)_{12}(YX)_{12}-\tr(XY)X_{12}Y_{12}
 =D_XY_{12}^2+D_YX_{12}^2.
\]
Consequently,
\begin{equation}\label{eq:ss}
 \boxed{\Del_{sz}H=D_XY_{12}^2+D_YX_{12}^2+2X_{12}Y_{12}\ge0.}
\end{equation}
For two different spokes in a wheel with positive rim weights, both segments contain a positive $R$ factor. Thus their upper-right entries are positive and the inequality is strict.

\subsection{Two inverted rim variables}
Conjugate both segments by $S$. Write
\[
 X'=\begin{pmatrix}a&b\\c&d\end{pmatrix},\qquad
 Y'=\begin{pmatrix}e&f\\g&h\end{pmatrix}.
\]
Since $SUS=\operatorname{diag}(1,0)$, direct differentiation and expansion yield
\begin{align}
 \Del_{tz}H
 &=D_XD_Y-adD_Y-ehD_X+2ae\label{eq:rr}\\
 &=D_XD_Y+D_Y(D_X-ad)+D_X(D_Y-eh)+2(ae-D_XD_Y)\notag\\
 &\ge D_XD_Y>0.\notag
\end{align}
Every parenthesized expression is nonnegative by Lemma~\ref{lem:invariants}. This treats any two different rim variables.

\subsection{One spoke and one inverted rim variable}
Here use the original, unconjugated entries
\[
 X=\begin{pmatrix}a&b\\c&d\end{pmatrix},\qquad
 Y=\begin{pmatrix}e&f\\g&h\end{pmatrix}.
\]
The relevant cyclic expression is $\tr((I+sE)X(I+tU)Y)-2$. Expanding its four coefficients gives
\begin{equation}\label{eq:mixed}
 K:=-\Del_{st}H=a(a-b)D_Y+D_Xh(f+h)-2a(f+h).
\end{equation}
Set $y=f+h$, $u=D_X(a-b)-a$ and $v=D_Yh-y$. Lemma~\ref{lem:invariants} gives $u,v\ge0$, and the following exact certificate proves the required sign:
\begin{equation}\label{eq:certificate}
 \boxed{D_XD_YK=(D_Ya-D_Xy)^2+D_Y^2a u+D_X^2y v\ge0.}
\end{equation}
Equivalently, division by $D_XD_Y$ and the inequalities in \eqref{eq:uv} bound $K/(D_XD_Y)$ below by $(a/D_X-y/D_Y)^2$.

For completeness, this sign is strict on a positive wheel. Mark $P_i$ and $R_j$, with the cyclic cut immediately before $P_i$. If $j\ne i$, the segment $X$ begins with $R_iP_{i+1}$. For positive parameters $I(R(t)P(s))=st^2>0$, and subsequent right multiplications preserve strict positivity. Hence $u>0$. If $j=i$, then $X=I$ and $Y$ contains all other complete $PR$ pairs. Building $Y$ from the right, its final pair already has $J(P(s)R(t))=st^2>0$; further left multiplications preserve strict positivity. Thus $v>0$ in this case. Since $a,y,D_X,D_Y$ are positive, \eqref{eq:certificate} is strict.

\subsection{Returning to the original weights}
Let $B=\prod_i r_i$. For distinct variables the monomial prefactor in \eqref{eq:transfer} contributes no mixed derivative to $\log Z$. The change $t_i=1/r_i$ has derivative $-r_i^{-2}$. Consequently,
\begin{align*}
 \Del_{s_i s_j}Z&=B^2\Del_{s_i s_j}H,\\
 \Del_{r_i r_j}Z&=\frac{B^2}{r_i^2r_j^2}\Del_{t_i t_j}H,\\
 \Del_{s_i r_j}Z&=-\frac{B^2}{r_j^2}\Del_{s_i t_j}H.
\end{align*}
The signs established in \eqref{eq:ss}, \eqref{eq:rr}, and \eqref{eq:certificate} prove Theorem~\ref{thm:main}.

\section{A sharp relaxation parameter and whirls}
The same certificates prove a slightly stronger classification. Let $\mathcal R$ be the complete set of rim edges and define, for $\lambda\ge0$,
\[
 Z_\lambda=Z_{W_n}+\lambda\prod_{i=1}^n r_i.
\]
This gives the formerly forbidden rim cycle a separate weight $\lambda$ while leaving every forest weight unchanged.

\begin{theorem}\label{thm:parameter}
Fix $n\ge3$. The polynomial $Z_\lambda$ is Rayleigh for every positive assignment of its edge variables if and only if $0\le\lambda\le1$.
\end{theorem}
\begin{proof}
Replace $H$ by $H_c=\tr\prod PR-c$, with $c=2-\lambda$. For $1\le c\le2$, the spoke--spoke difference is
\[
 D_XY_{12}^2+D_YX_{12}^2+cX_{12}Y_{12}\ge0.
\]
In the conjugated coordinates the rim--rim difference is
\[
 D_XD_Y-adD_Y-ehD_X+cae\ge(c-1)D_XD_Y\ge0.
\]
For a mixed pair, using the notation of \eqref{eq:mixed},
\[
 -\Del_{st}H_c=K+(2-c)a(f+h)\ge0.
\]
The same inverse-variable signs finish the sufficient direction.

For necessity suppose $\lambda>1$ and temporarily set every spoke weight to zero. Then
\[
 Z_\lambda=\prod_i(1+r_i)+(\lambda-1)\prod_i r_i.
\]
For two different rim edges $j,k$, put $A=\prod_{i\ne j,k}(1+r_i)$ and $B=(\lambda-1)\prod_{i\ne j,k}r_i$. The resulting polynomial in their weights is $A(1+r_j)(1+r_k)+Br_jr_k$, so its Rayleigh difference is $-AB<0$. Continuity preserves this negative sign when all spoke weights are sufficiently small and positive. This excludes every $\lambda>1$.
\end{proof}

\begin{corollary}
The independent-set polynomial of every rank-$n$ whirl, $n\ge3$, is Rayleigh for all positive weights.
\end{corollary}
\begin{proof}
The whirl is obtained by relaxing the rim circuit-hyperplane of the wheel matroid. Its bases are the spanning trees of the wheel together with $\mathcal R$ \cite{DSW}. All proper subsets of $\mathcal R$ were already independent, so its independent-set polynomial is exactly $Z_1$. Apply Theorem~\ref{thm:parameter}.
\end{proof}
The algebraic parameter classification is established here. Priority for its matroid interpretation has not been exhaustively checked.

\section{Consequences with precise scope}
\begin{corollary}\label{cor:closure}
Every minor of a wheel, and every graph obtained by repeated graph one-sums and two-sums of wheels and series-parallel graphs, has a Rayleigh forest polynomial at arbitrary positive weights. The same holds after any feasible specification that selected edges are present and other selected edges are absent.
\end{corollary}
\begin{proof}
Deletion is specialization at weight zero; contraction is the leading coefficient as that weight tends to infinity. Rayleigh inequalities pass to both limits. Feasible conditioning gives precisely these operations, with loops discarded.

One-sums multiply forest polynomials. For a standard two-sum along a non-loop, non-bridge edge, write $A=A_0+zA_1$ and $B=B_0+zB_1$. The resulting polynomial is
\[
 Z=A_0B_1+A_1B_0-A_1B_1
 =B_1A\!\left(\frac{B_0-B_1}{B_1}\right).
\]
Here $B_0-B_1\ge0$: the forests extendible by the gluing edge form a subset of the forests avoiding that edge. For two edges on the $A$ side,
\[
 \Del_{ef}Z=B_1^2\Del_{ef}A\!\left(\frac{B_0-B_1}{B_1}\right).
\]
For one edge on each side, direct expansion gives
\[
 \Del_{ef}Z=\Del_{ez}A\;\Del_{zf}B.
\]
These standard closure identities \cite{Wagner,Erickson} propagate the wheel theorem. Series-parallel graphs satisfy the weighted inequality \cite{Erickson}.
\end{proof}

\begin{corollary}[Cylinder and concentration inequalities]
For any law in Corollary~\ref{cor:closure}, put $p_e=\Pr(e\in A)$ and $N=|A|$. For every set $T$ of edges,
\[
 \Pr(T\subseteq A)\le\prod_{e\in T}p_e,\qquad
 \Pr(T\cap A=\varnothing)\le\prod_{e\in T}(1-p_e).
\]
For every real $\theta$,
\begin{equation}\label{eq:mgf}
 \mathbb E e^{\theta N}\le\prod_e(1-p_e+p_e e^\theta).
\end{equation}
Thus the usual exponential Markov bounds for the sum of independent Bernoulli variables with these same marginals also bound the upper and lower tails of $N$.
\end{corollary}
\begin{proof}
Under every feasible partial conditioning, pairwise negative correlation says that conditioning another edge to be present cannot increase a remaining edge's inclusion probability. Apply this repeatedly to the chain rule. For simultaneous absence use the indicators $1-X_e$, whose pairwise covariances have the same signs, and deletion conditioning. Impossible events have probability zero and need no conditioning.

For $\theta\ge0$, expand $\prod_e[1+(e^\theta-1)X_e]$ and apply the inclusion inequality termwise. For $\theta<0$, use
\[
 e^{\theta N}=e^{\theta |E|}\prod_e[1+(e^{-\theta}-1)(1-X_e)]
\]
and the absence inequality. This proves \eqref{eq:mgf}.
\end{proof}

\begin{corollary}[Weighted connected spanning subgraphs]
On every wheel, the law on connected spanning edge sets $A$ with weight proportional to $\prod_{e\in A}w_e$ has strictly negative correlation for each pair of distinct edges, for all positive weights.
\end{corollary}
\begin{proof}
In a connected plane graph $G$, an edge set $A$ is spanning and connected exactly when its complementary dual edge set is a forest of $G^*$. Its weight transforms by
\[
 \prod_{e\in A}w_e=\left(\prod_{e\in E}w_e\right)
 \prod_{e^*\notin A}w_e^{-1}.
\]
The plane dual of a wheel is another wheel. Complementing both indicators preserves their covariance, so Theorem~\ref{thm:main} applies.
\end{proof}

\section{Verification and remaining boundary}
The proof is algebraic for every $n$, not an extrapolation from computation. Separate symbolic checks verified the three generic two-by-two identities. Direct acyclic-subset enumeration agreed coefficient by coefficient with the transfer polynomial for $n=3,4,5,6$. A separate integer-weight check for $n=3,\ldots,8$ matched the partition function, all first derivatives, and all mixed second derivatives against enumeration, and found strict positivity for all 365 edge pairs tested. These finite checks detect transcription errors; the lemmas establish the infinite statement.

The argument proves pairwise Rayleigh inequalities and the stated consequences. It does not prove negative association for arbitrary increasing events, real stability of the unrooted forest polynomial, the random-cluster conjecture for every $0<q<1$, or the forest conjecture for all finite graphs. The reusable matrix certificate covers cyclic products of the displayed two matrix types; an arbitrary graph does not automatically admit that representation.

\begin{thebibliography}{9}
\bibitem{Bado} I. O. Bado, \emph{Negative edge correlation for uniform random forests on wheel graphs}, Statistics \& Probability Letters, available online 2026, PII S016771522600324X. \url{https://www.sciencedirect.com/science/article/abs/pii/S016771522600324X}.
\bibitem{CR} R. A. \c{C}i\c{c}eksiz and M. Ravichandran, \emph{Negative Correlation for Forests and the $q<1$ Random-Cluster Model}, arXiv:2610.00549v1 (2026). \url{https://arxiv.org/abs/2610.00549}.
\bibitem{Wagner} D. G. Wagner, \emph{Negatively Correlated Random Variables and Mason's Conjecture}, arXiv:math/0602648 (2006), especially Lemma 4.1 and Theorem 5.8. \url{https://arxiv.org/abs/math/0602648}.
\bibitem{Erickson} A. Erickson, \emph{Sums of squares and negative correlation for spanning forests of series parallel graphs}, Australasian Journal of Combinatorics \textbf{52} (2012), 75--87. \url{https://ajc.maths.uq.edu.au/pdf/52/ajc_v52_p075.pdf}.
\bibitem{DSW} G. Denham, M. Schulze and U. Walther, \emph{Matroid connectivity and singularities of configuration hypersurfaces}, Letters in Mathematical Physics \textbf{111} (2021), article 11. \url{https://doi.org/10.1007/s11005-020-01352-3}.
\bibitem{VMRC} QEEDDD, \emph{Verifiable Mathematical Reasoning Corpus}, research drafts [016] and [160], repository revision \texttt{38348353416f1f7fc50cecb93c3b554a00260c13}. \url{https://github.com/QEEDDD/Verifiable-Mathematical-Reasoning-Corpus}.
\end{thebibliography}
\end{document}
